%!TEX ROOT=../zpevnik.tex

[Em]Kdo [G]se [Am]večer [Em]hájem [Am]vrací,[F] ten ať [Em]klopí[G7] zra[C]ky, [Em]
ať[G] je [Am]nikdy[Em] neo[Am]brací [Dm]k [G7]vrbě [Em]křivola[E]ký. [A] [C] [Em]
Jin[G]ak [Am]jeho [Em]oči [Am]zjistí, [F]i když [Em]se to [G7]nez[C]dá, [Em]
že[G] na [Am]vrbě [Em]kromě [Am]listí, [Dm]visí [F]malá hvěz[Am]da.

(R): Vidě[C]li jsme jednou v [F]lukách,
plakat [C]na tý vrbě [F]kluka,
který [Dm]pevně věřil [B]tomu,
že ji [D7]sundá z toho [G]stromu.

Kdo o hvězdy jeví zájem, zem když večer chladne,
ať jde klidně přesto hájem, hvězda někdy spadne.
Ať se pro ni rosou brodí a pak vrbu najde si,
a pro ty co kolem chodí na tu větev zavěsí.

(R):

